\begin{description}
   \item[1行目] データファイル\verb|baseballDecade.csv|を読み込みます。
   \item[2-8行目] 分析用のデータに加工しています。
      \begin{description}
         \item[3行目] フィルターをかけてデータを野手だけのものにします。
         \item[4行目] フィルターをさらにかけて，50試合場出ている選手に限定します。
         \item[5行目] チームがTigersの選手だけに絞っています。
         \item[6行目] 年俸のデータは歪んでいるので，対数をとって正規分布の形に近似させました\footnote{データを変形させずに確率分布の方で対応するのがよい，と常々話をしてきましたが，ここではデータを変形させています。実は左に歪んだような分布は対数正規分布というのがあり，そこからのデータ生成を考えても良いのですが，対数正規分布はデータの対数が従う正規分布というだけなので，データの方を変えてしまった方が説明しやすいのです。比率や度数など，データの生成メカニズムがそもそも違うばあいは，データを変換せずに使うのが正しいやり方です。}。
         \item[7行目] 使う変数だけに絞っています。
         \item[8行目] 行番号を\verb|ID|という変数名にして追加しました
      \end{description}
